\scanpage{169}

sous-groupe distingué ds un groupe ambiant, la valeur du groupe quotient. Les deux carrés de ce diagramme sont cartésiens et cocartésiens -- en particulier, $B$ avec son sous-groupe invariant \deleted{\unsure{$g$}} $\widetilde B''$, se déduit de $(B'\supset\widetilde U)$ et de l'action de $B'$ sur $\widetilde B''$ (induite par (33)), de la $=$ façon que $(G\supset G_0)$ se déduit de $(B'\supset\widetilde U)$ et de l'action (33) de $B'$ sur $G_0$.

Le groupe $\widetilde B''$ \deleted{\unsure{est, ici}} sera, dans les cas qui nous occupent, ``petits'' et bien \unsure{pigeables|pignables|figurables} \inserted{(en plus d'être plongé dans $G_0$)}, au $=$ titre que $B'$ auquel il ressemble comme un frère, et il en sera de $=$ de $B$.

Nous avons exprimé pour l'instant \inserted{dans (26 c))} la donnée d'un sous-groupe $Z$ de $\unsure{\mathfrak S|\mathscr S}_S$ commutant à l'action de $G$, mais il reste encore à exprimer que $Z$ est donné en \deleted{tant} que sous-groupe central de $G$. Cela s'explicite par la donnée d'un homomorphisme \deleted{$\unsure{Z}$}

\[ Z\longrightarrow G \tag*{(40)} \]

tel que les deux actions de $Z$ sur $S$ -- celles du début, et celle déduite de (40), soient égales.

On voit alors tout de suite que cela signifie que \deleted{l'on a} \inserted{(40) provient} d'un \uline{scindage} \deleted{\unsure{central}} \inserted{``central''} de l'extension $B$ de $Z$ par $B'$, \deleted{qui est}

\diagram{\begin{tikzcd}[ampersand replacement=\&,column sep=small,row sep=tiny]
(41)\quad\left\{\vphantom{\begin{matrix}X\\X\\X\end{matrix}}\right.\ 0\arrow[r,dash]\&B'\arrow[r]\&B\arrow[r,dash]\&Z\arrow[r,dash]\arrow[l,bend left=38]\&0
\end{tikzcd}}

définissant \inserted{donc} un isomorphisme $B\simeq\deleted{\unsure{B'}}Z\times B'$ avec en plus $Z$\inserted{$\subset\unsure{B|B'}$} opérant \uline{trivialement} sur $G_0$ (donc sur $G$)

% archive sheet number: 168
